\section{Exact separator audit}
\label{app:separator}

This appendix records the algebraic boundary behind
\eqref{eq:separator-c4}.  Starting from the cumulants
\(\ell_2,\ldots,\ell_6\), the verifier reconstructs the central moments
\[
 m_0=1,\quad m_1=0,\quad m_2=\ell_2,\quad m_3=\ell_3,
\]
\[
 m_4=\ell_4+3\ell_2^2,\qquad
 m_5=\ell_5+10\ell_3\ell_2,
\]
\[
 m_6=\ell_6+15\ell_4\ell_2+10\ell_3^2+15\ell_2^3,
\]
and then forms \(\det[m_{i+j}]_{i,j=0}^3\) directly.  This independently
regenerates the thirteen-term \(\Cfour\) polynomial; it does not import the
expansion used for the Riemann kernel.

The recurrence \eqref{eq:separator-N-recurrence} is then applied over
\(\mathbb Q[t]\).  Substitution of \eqref{eq:separator-logjets} shows that
each of the thirteen determinant monomials has common denominator \(P^{12}\).
After collecting the numerator \(H\), exact rational arithmetic verifies
\[
 \deg H=72,\qquad [t^j]H=[t^{72-j}]H>0\quad(0\leq j\leq72),
\]
with
\[
 \min_{0\leq j\leq72}[t^j]H=\frac3{65536}.
\]
For a printable certificate, write
\[
 2^{54}H(t)=\sum_{j=0}^{72}a_jt^j,
 \qquad a_j=a_{72-j}.
\]
The 37 independent integer coefficients are:
\begin{center}
\tiny
\begin{tabular}{@{}r r@{\qquad}r r@{}}
\toprule
$j$ & $a_j$ & $j$ & $a_j$ \\
\midrule
0 & $824633720832$ & 19 & $6987741510406573876435979520$ \\
1 & $48712961228800$ & 20 & $17092065810580803244866465328$ \\
2 & $1429432562155520$ & 21 & $39268975295924090460974938880$ \\
3 & $27770895893217280$ & 22 & $84898287500864044148851126266$ \\
4 & $401708065273655296$ & 23 & $173000552553192254015824625680$ \\
5 & $4613079582269972480$ & 24 & $332750476261206454685845998080$ \\
6 & $43792294989734840832$ & 25 & $604872001890033837007447191280$ \\
7 & $353349130016052428800$ & 26 & $1040322081625845770286033564593$ \\
8 & $2472885380960577236992$ & 27 & $1694567248719256234936725175440$ \\
9 & $15243207814279730964480$ & 28 & $2616435044092694236816009466532$ \\
10 & $83763600430517403650816$ & 29 & $3832176384214979954574258944560$ \\
11 & $414327966630038761815040$ & 30 & $5327751824843881176247055457873$ \\
12 & $1859451652398640291313152$ & 31 & $7034632066783174173709968967280$ \\
13 & $7621494792057501544629760$ & 32 & $8825381090424751586774701723088$ \\
14 & $28689460125397617502291392$ & 33 & $10523927205327676216804189288720$ \\
15 & $99652010444258917615558400$ & 34 & $11931512643001413047042570621636$ \\
16 & $320698275507551388688130816$ & 35 & $12863883307375841645695614141760$ \\
17 & $959582391519228973730931840$ & 36 & $13190420295966739733605724579768$ \\
18 & $2677774483831421514536291072$ & & \\
\bottomrule
\end{tabular}
\end{center}

This table is generated by
\path{scripts/generate_separator_coefficients.py}. Its machine-readable
companion \path{generated/separator-coefficients.json} has SHA-256 prefix
\texttt{87af00802c0929ca}; the complete hash is recorded in the release
manifest.
Coefficient positivity is a proof on the full domain \(t>0\), not a sampled
test. Exact spot reconstructions at \(t=1/3,1,4\) independently compare the
central determinant with \(H/P^{12}\).

For the Fourier step, the identities
\[
 w+w^{-1}=u,\quad w^2+w^{-2}=u^2-2,\quad
 w^3+w^{-3}=u^3-3u
\]
give \eqref{eq:separator-cubic} directly. Put
\[
 r=e^{1/64},\qquad x=r^2=e^{1/32}.
\]
Exact expansion gives
\[
 \operatorname{disc}(R_c)=4r^{10}G(x),
\]
where
\begin{align*}
G(x)={}&432x^{13}-4968x^9+900x^8+16200x^6+19044x^5\\
&-72600x^4+20000x^3+62100x^2-54334x+13225.
\end{align*}
Writing \(y=x-1\) gives
\begin{align*}
G(1+y)={}&-1-10y-12y^2+680y^3+11532y^4+96660y^5\\
&+365400y^6+569664y^7+512172y^8+303912y^9\\
&+123552y^{10}+33696y^{11}+5616y^{12}+432y^{13}.
\end{align*}
Moreover \(0<y<1/30\). Indeed,
\[
 e<1+1+\frac12+\frac16+\frac1{18}
 =\frac{49}{18}<\frac{11}{4},
\]
where \(1/18=(1/24)\sum_{j\geq0}4^{-j}\) bounds the tail beginning with
\(1/4!\). The binomial theorem gives
\[
 \frac{11}{4}<1+\frac{32}{30}
 +\frac{\binom{32}{2}}{30^2}+\frac{\binom{32}{3}}{30^3}
 =\frac{1891}{675}<\left(\frac{31}{30}\right)^{32}.
\]
The sum of the positive terms at \(y=1/30\) is exactly
\[
 \sum_{j=3}^{13}\frac{[y^j]G(1+y)}{30^j}
 =\frac{1621193195302591}{36905625000000000}<1.
\]
The two omitted terms \(-10y-12y^2\) are negative. Consequently
\[
 G(1+y)<-1+
 \sum_{j=3}^{13}\frac{[y^j]G(1+y)}{30^j}<0,
\]
and hence \(\operatorname{disc}(R_c)<0\) without numerical evaluation.
